% Appendix for ProgressProbe paper
% Contains supplementary tables, ablation details, and error analysis breakdowns.

\appendix

\section{Supplementary Material}
\label{sec:appendix}

%-------------------------------------------------------------------
\subsection{Encoder Ablation Complete Results}
\label{sec:app_encoder}

Table~\ref{tab:encoder_ablation_full} reports ProgressProbe performance across three text encoders with multiple random seeds.
All models use the same GRU architecture (2-layer, hidden$=$256, dropout$=$0.1) and training procedure (lr$=$1e-3, batch$=$32, early stopping with patience~5).
Evaluation is on the SWE-Agent conf${\geq}$2 test set ($N{=}1{,}824$).

\begin{table}[h]
\centering
\small
\begin{tabular}{llcccc}
\toprule
\textbf{Encoder} & \textbf{Seed} & \textbf{Macro-F1} & \textbf{F1\textsubscript{stag}} & \textbf{F1\textsubscript{prog}} & \textbf{Mean$\pm$Std} \\
\midrule
\multirow{2}{*}{SFR-Embedding (0.6B)} & 42 & 0.9055 & 0.8593 & 0.9518 & \multirow{2}{*}{0.9185$\pm$0.0113} \\
                                       & 123, 456 & \multicolumn{3}{c}{\emph{(per-seed detail not recorded)}} & \\
\midrule
\multirow{4}{*}{E5-large-v2 (335M)}   & 42  & 0.9363 & 0.9067 & 0.9660 & \multirow{4}{*}{0.9381$\pm$0.0020} \\
                                       & 123 & 0.9403 & --- & --- & \\
                                       & 456 & 0.9378 & --- & --- & \\
                                       & \emph{mean} & --- & 0.9173$\pm$0.010 & 0.9590$\pm$0.006 & \\
\midrule
\multirow{4}{*}{BGE-M3 (568M)}        & 42  & 0.9413 & 0.9147 & 0.9679 & \multirow{4}{*}{0.9084$\pm$0.0529} \\
                                       & 123 & 0.8474 & --- & --- & \\
                                       & 456 & 0.9364 & --- & --- & \\
                                       & \emph{mean} & --- & 0.8834$\pm$0.057 & 0.9333$\pm$0.049 & \\
\bottomrule
\end{tabular}
\caption{Encoder ablation with multi-seed evaluation. BGE-M3 seed~123 is an outlier (0.847) due to split-by-instance variation in the train/test partition; the seed~42 result (0.941) is used as the primary comparison in the main paper.}
\label{tab:encoder_ablation_full}
\end{table}

All three encoders exceed 0.90 macro-F1 at seed~42, with E5-large-v2 and BGE-M3 surpassing the default SFR-Embedding by 3.1 and 3.6\,pp respectively.
However, BGE-M3 exhibits high variance across seeds ($\sigma{=}0.053$), driven by the seed~123 outlier, making E5-large-v2 the most stable alternative ($\sigma{=}0.002$).

%-------------------------------------------------------------------
\subsection{Window Size Ablation}
\label{sec:app_window}

Table~\ref{tab:window_ablation} reports ProgressProbe performance as a function of window size $k \in \{3, 5, 7\}$.
All runs use the SFR-Embedding encoder with seed~42.

\begin{table}[h]
\centering
\small
\begin{tabular}{ccccc}
\toprule
\textbf{Window $k$} & \textbf{Macro-F1} & \textbf{F1\textsubscript{stag}} & \textbf{F1\textsubscript{prog}} & \textbf{$\Delta$ vs.\ $k{=}5$} \\
\midrule
3 & 0.8721 & 0.7784 & 0.9659 & $-$5.9\,pp \\
5 & \textbf{0.9309} & \textbf{0.8917} & \textbf{0.9701} & --- \\
7 & 0.9133 & 0.8791 & 0.9475 & $-$1.8\,pp \\
\bottomrule
\end{tabular}
\caption{Window size ablation. $k{=}5$ is optimal. Smaller windows ($k{=}3$) lack sufficient context for stagnation detection ($-$5.9\,pp, concentrated in F1\textsubscript{stag}). Larger windows ($k{=}7$) introduce noise that dilutes stagnation signals ($-$1.8\,pp).}
\label{tab:window_ablation}
\end{table}

The performance drop at $k{=}3$ is concentrated in stagnation detection (F1\textsubscript{stag} drops 11.3\,pp), confirming that stagnation patterns require multi-step context to distinguish from single noisy actions.
The milder degradation at $k{=}7$ ($-$2.3\,pp on F1\textsubscript{prog}) suggests that longer windows dilute the temporal signal by mixing progress and stagnation steps.

%-------------------------------------------------------------------
\subsection{LF Temporal Dynamics Statistics}
\label{sec:app_temporal}

Table~\ref{tab:lf_lag} reports the lag distribution of each labeling function relative to stagnation onset, measured over all trajectories in the SWE-Agent dataset.
Negative lag indicates the LF fires \emph{before} the onset boundary (leading); positive lag indicates firing \emph{after} (lagging).

\begin{table}[h]
\centering
\small
\begin{tabular}{lccc}
\toprule
\textbf{Labeling Function} & \textbf{Mean Lag (steps)} & \textbf{\% Leading} & \textbf{\% Lagging} \\
\midrule
Command repetition ($\lambda_1$) & $-$1.87 & 36.1\% & --- \\
File-diff delta ($\lambda_2$) & $-$1.20 & 17.1\% & --- \\
Outcome backprop.\ ($\lambda_3$) & $+$6.70 & --- & 48.4\% \\
\bottomrule
\end{tabular}
\caption{Temporal lag statistics for each labeling function. $\lambda_1$ and $\lambda_2$ are leading indicators of stagnation; $\lambda_3$ is a lagging indicator due to its retrospective outcome-based design.}
\label{tab:lf_lag}
\end{table}

The Kruskal--Wallis test confirms that lag distributions differ significantly across LFs ($H{=}274.3$, $p{=}2.74 \times 10^{-60}$).
All pairwise comparisons are significant after Bonferroni correction (Table~\ref{tab:lf_pairwise}).

\begin{table}[h]
\centering
\small
\begin{tabular}{lc}
\toprule
\textbf{Comparison} & \textbf{$p$-value (Bonferroni)} \\
\midrule
$\lambda_1$ vs.\ $\lambda_2$ (cmd\_rep vs.\ file\_diff) & $1.44 \times 10^{-26}$ \\
$\lambda_1$ vs.\ $\lambda_3$ (cmd\_rep vs.\ outcome\_bp) & $6.67 \times 10^{-39}$ \\
$\lambda_2$ vs.\ $\lambda_3$ (file\_diff vs.\ outcome\_bp) & $3.47 \times 10^{-47}$ \\
\bottomrule
\end{tabular}
\caption{Pairwise Dunn's post-hoc tests with Bonferroni correction. All three LFs have statistically distinct temporal profiles.}
\label{tab:lf_pairwise}
\end{table}

%-------------------------------------------------------------------
\subsection{Structural Feature Importance}
\label{sec:app_features}

Table~\ref{tab:feature_list} lists the 30 hand-crafted structural features used by the non-NLP baselines (RF, GB, LR).
Features are grouped by category.

\begin{table}[h]
\centering
\small
\begin{tabular}{lll}
\toprule
\textbf{Category} & \textbf{Feature} & \textbf{Description} \\
\midrule
\multirow{4}{*}{Window stats}
  & \texttt{num\_steps} & Total steps in window \\
  & \texttt{num\_actions} & Agent action count \\
  & \texttt{num\_observations} & Environment observation count \\
  & \texttt{action\_ratio} & Actions / total steps \\
\midrule
\multirow{7}{*}{Text length}
  & \texttt{mean\_text\_len} & Mean text length across steps \\
  & \texttt{std\_text\_len} & Std of text lengths \\
  & \texttt{max\_text\_len} & Maximum text length \\
  & \texttt{min\_text\_len} & Minimum text length \\
  & \texttt{mean\_action\_len} & Mean action text length \\
  & \texttt{mean\_obs\_len} & Mean observation text length \\
  & \texttt{text\_len\_trend} & Linear slope of text length over steps \\
\midrule
\multirow{4}{*}{Command stats}
  & \texttt{num\_commands} & Extracted command count \\
  & \texttt{mean\_cmd\_len} & Mean command length \\
  & \texttt{max\_cmd\_len} & Maximum command length \\
  & \texttt{unique\_cmd\_ratio} & Unique commands / total \\
\midrule
\multirow{1}{*}{Repetition}
  & \texttt{max\_repeat\_ratio} & Most-repeated cmd / total \\
\midrule
\multirow{6}{*}{Command types}
  & \texttt{num\_edit\_cmds} & Edit command count \\
  & \texttt{num\_search\_cmds} & Search command count \\
  & \texttt{num\_nav\_cmds} & Navigation command count \\
  & \texttt{num\_submit\_cmds} & Submit command count \\
  & \texttt{edit\_ratio} & Edit commands / total \\
  & \texttt{search\_ratio} & Search commands / total \\
\midrule
\multirow{3}{*}{Cmd similarity}
  & \texttt{avg\_adjacent\_cmd\_sim} & Mean adjacent command similarity \\
  & \texttt{max\_adjacent\_cmd\_sim} & Max adjacent command similarity \\
  & \texttt{max\_consecutive\_similar} & Longest run of similar commands \\
\midrule
\multirow{4}{*}{Error signals}
  & \texttt{num\_errors} & Tracebacks / error strings \\
  & \texttt{num\_test\_outputs} & Test PASSED/FAILED outputs \\
  & \texttt{num\_file\_displays} & File display markers \\
  & \texttt{num\_nonzero\_exits} & Non-zero exit codes \\
\midrule
\multirow{1}{*}{Metadata}
  & \texttt{is\_short\_trajectory} & Binary: trajectory $<$ threshold \\
\bottomrule
\end{tabular}
\caption{Complete list of 30 structural features. Feature importance ranking from RF/GB models is pending a dedicated analysis run; the feature definitions are provided here for reproducibility.}
\label{tab:feature_list}
\end{table}

%-------------------------------------------------------------------
\subsection{Data Scaling Complete Curves}
\label{sec:app_scaling}

Table~\ref{tab:data_scaling} reports ProgressProbe performance as a function of training data fraction on the SWE-Agent dataset.
All runs use the SFR-Embedding encoder with seed~42.

\begin{table}[h]
\centering
\small
\begin{tabular}{rcccc}
\toprule
\textbf{Data \%} & \textbf{Macro-F1} & \textbf{F1\textsubscript{stag}} & \textbf{F1\textsubscript{prog}} & \textbf{Accuracy} \\
\midrule
10\% & 0.8054 & 0.7108 & 0.9000 & 0.8514 \\
25\% & 0.8203 & 0.7316 & 0.9090 & 0.8640 \\
50\% & 0.8715 & 0.8078 & 0.9351 & 0.9030 \\
100\% & \textbf{0.9055} & \textbf{0.8593} & \textbf{0.9518} & \textbf{0.9282} \\
\bottomrule
\end{tabular}
\caption{ProgressProbe data scaling on embedding features. Performance improves monotonically with data size; 10\% data still achieves 0.81 macro-F1, with stagnation F1 as the primary beneficiary of additional data ($+$14.9\,pp from 10\% to 100\%).}
\label{tab:data_scaling}
\end{table}

\paragraph{Structural baseline data scaling (circular evaluation caveat).}
For completeness, we also evaluated the structural RF/GB baselines under the same data fractions (Table~\ref{tab:structural_scaling}).
\textbf{These numbers must be interpreted with extreme caution}: structural features are computed from the same signals that the labeling functions use to produce silver labels (\S\ref{sec:circular}).
The near-perfect performance at 10\% data (${\geq}0.987$ macro-F1) is a \emph{circular evaluation artifact}---the features trivially reverse-engineer the LFs regardless of training set size---and \textbf{should not be cited as evidence of data-scarcity robustness}.

\begin{table}[h]
\centering
\small
\begin{tabular}{rcc}
\toprule
\textbf{Data \%} & \textbf{RF Macro-F1} & \textbf{GB Macro-F1} \\
\midrule
10\% & 0.9909 & 0.9873 \\
25\% & 0.9916 & 0.9930 \\
50\% & 0.9951 & 0.9951 \\
100\% & 0.9951 & 0.9951 \\
\bottomrule
\end{tabular}
\caption{Structural baseline data scaling. $^\dagger$\textbf{Circular evaluation}: structural features $\approx$ LF inputs, so near-perfect performance at all data fractions reflects LF reconstructibility, not genuine detection accuracy. Included for completeness only.}
\label{tab:structural_scaling}
\end{table}

%-------------------------------------------------------------------
\subsection{Error Analysis Supplementary Tables}
\label{sec:app_errors}

Tables~\ref{tab:error_summary}--\ref{tab:error_patterns} provide the full stratified error analysis of ProgressProbe (GRU, seed~42) on the SWE-Agent conf${\geq}$2 test set.

\begin{table}[h]
\centering
\small
\begin{tabular}{lrrr}
\toprule
\textbf{Stratum} & \textbf{$N$} & \textbf{Errors} & \textbf{Error Rate} \\
\midrule
\multicolumn{4}{l}{\emph{Overall}} \\
\quad Total & 1,824 & 131 & 7.18\% \\
\quad False Positives (pred stag, true prog) & --- & 85 & 64.9\% of errors \\
\quad False Negatives (pred prog, true stag) & --- & 46 & 35.1\% of errors \\
\midrule
\multicolumn{4}{l}{\emph{By LF confidence}} \\
\quad conf$=$2 (2/3 LFs agree) & 1,780 & 126 & 7.08\% \\
\quad conf$=$3 (unanimous) & 44 & 5 & 11.36\%$^*$ \\
\midrule
\multicolumn{4}{l}{\emph{By task outcome}} \\
\quad Resolved (task succeeded) & 77 & 1 & 1.30\% \\
\quad Unresolved (task failed) & 1,747 & 130 & 7.44\% \\
\midrule
\multicolumn{4}{l}{\emph{By trajectory position}} \\
\quad Early (0--33\%) & 351 & 23 & 6.55\% \\
\quad Middle (33--66\%) & 353 & 36 & 10.20\% \\
\quad Late (66--100\%) & 1,070 & 69 & 6.45\% \\
\bottomrule
\end{tabular}
\caption{Error stratification by confidence, outcome, and position. $^*$conf$=$3 error rate is elevated but based on only 44 samples; the difference is not statistically reliable. Middle-trajectory windows have the highest error rate (10.2\%), consistent with the exploration--stagnation boundary being most ambiguous at trajectory midpoints.}
\label{tab:error_summary}
\end{table}

\begin{table}[h]
\centering
\small
\begin{tabular}{lc}
\toprule
\textbf{Metric} & \textbf{Value} \\
\midrule
\multicolumn{2}{l}{\emph{LF conflict analysis}} \\
\quad Base rate (all test windows) & 31.0\% \\
\quad Error windows with LF conflict & 39.7\% \\
\midrule
\multicolumn{2}{l}{\emph{False positive (FP) patterns}} \\
\quad Mean predicted prob (stagnation) & 0.922 \\
\quad FP with prob $>$ 0.9 & 73\% \\
\midrule
\multicolumn{2}{l}{\emph{Error concentration}} \\
\quad Top 2 trajectories share of errors & 18.3\% \\
\bottomrule
\end{tabular}
\caption{Error characteristics. FPs are highly confident (mean stagnation prob$=$0.922), indicating systematic blind spots rather than borderline cases. LF conflict is a strong error predictor (40\% vs.\ 31\% base rate). Errors are non-uniformly distributed, with the top~2 trajectories contributing 18.3\% of all misclassifications.}
\label{tab:error_patterns}
\end{table}
